\section{Nuclear averaging and weighted matrix blocks}\label{r:sec:weighted}

In this section \(A\subseteq D\) is a unital inclusion, \(A\) is nuclear,
and \(D\) has real rank zero and nonempty tracial state space. For a
prescribed finite set of unitary tests, nuclear averaging turns a
projection into a nearly central positive contraction. We will approximate
this contraction in trace by a weighted sum of projections, keeping their
total weighted squared commutator norm small. Each projection will also
be represented as a matrix block over a projection \(q\), with each
compressed test close to a scalar matrix over \(q\). This lets us sample
a subprojection \(e\leq q\) simultaneously in the matrix coordinates
while retaining trace and commutator estimates, without requiring \(e\)
to be central. We first choose the averaging map, independently of its
input projection.

\begin{lemma}[Averaging with trace control]\label{r:weighted:diagonal}
For a finite \(K\subseteq\U(A)\) and \(\mu>0\), there are contractions
\(v_1,\ldots,v_m\in A\) and numbers \(\lambda_l\geq0\) with
\(\sum_l\lambda_l=1\) such that the map
\[
 \Phi(b)=\sum_{l=1}^m\lambda_l v_lbv_l^*,\qquad b\in D,
\]
satisfies, for every projection \(q\in D\), every \(x\in K\), and every
\(\tau\in\T(D)\),
\begin{align}
 0\leq\Phi(q)&\leq1,\qquad
 \tau(q)-\mu\leq\tau(\Phi(q))\leq\tau(q),\label{r:eq:diagonal-mass}\\
 \normone{\Phi(q)-x\Phi(q)x^*}{\tau}
 &\leq\mu\tau(q).\label{r:eq:diagonal-central}
\end{align}
\end{lemma}

\begin{proof}
Let \(E=A\widehat\otimes A\) denote the Banach projective tensor product,
with norm \(\norm{\cdot}_\pi\), multiplication map
\(m_A(a\otimes b)=ab\), and bimodule actions
\(x\cdot(a\otimes b)=xa\otimes b\) and
\((a\otimes b)\cdot x=a\otimes bx\).
Haagerup's virtual diagonal theorem
\cite[Theorem~3.1]{Haa83} gives a virtual diagonal \(M\in E^{**}\) in
the weak-star closed convex hull of
\(\{v\otimes v^*:v\in A,\ \norm v\leq1\}\).
Its bimodule commutators vanish and \(m_A^{**}(M)=1\).
Apply the second adjoint of the bounded linear map
\[
 T:E\longrightarrow E^K\oplus A,\qquad
 T(d)=\big((x\cdot d-d\cdot x)_{x\in K},m_A(d)\big).
\]
Its value at \(M\) is \((0,\ldots,0,1)\), in the canonical copy of
\(E^K\oplus A\). Weak-star convergence in the bidual restricts to weak
convergence in this copy, so this point lies in the weak closure of the
image of the finite convex hull. A convex subset of a Banach space has the
same weak and norm closures, by Hahn--Banach separation. We may therefore choose a finite
convex combination \(d=\sum_l\lambda_l v_l\otimes v_l^*\) such that
\begin{equation}\label{r:eq:approx-diagonal}
 \norm{x\cdot d-d\cdot x}_\pi<\mu\quad(x\in K),\qquad
 \norm{1-\sum_l\lambda_l v_lv_l^*}<\mu.
\end{equation}

Both sums \(\sum_l\lambda_l v_lv_l^*\) and
\(\sum_l\lambda_l v_l^*v_l\) are positive contractions. By traciality,
the positive element \(c=1-\sum_l\lambda_l v_l^*v_l\) satisfies
\(0\leq\tau(c)<\mu\). Since
\(\tau(q)-\tau(\Phi(q))=\tau(qc)\), this proves
\eqref{r:eq:diagonal-mass}.
For the centrality estimate, the map
\(a\otimes b\mapsto aqb\) has norm at most \(\tau(q)\) from the
projective tensor norm to the tracial 1-seminorm: indeed,
\[
 \normone{aqb}{\tau}
 \leq\normtwo{aq}{\tau}\normtwo{qb}{\tau}
 \leq\norm a\norm b\,\tau(q).
\]
Apply this map to the first inequality in \eqref{r:eq:approx-diagonal}
and multiply by the unitary \(x^*\).
\end{proof}

The following tracial square-root estimate is the form of the
Powers--St\o rmer inequality needed below. Their original
\cite[Lemma~4.1]{PowersStormer1970} concerns positive Hilbert-space
operators with the ordinary trace; the short proof here gives the
arbitrary tracial $C^*$-algebra version directly.

\begin{lemma}[A square-root trace estimate]\label{r:weighted:square-root}
For positive \(b,c\in D\) and \(\tau\in\T(D)\),
\[
 \normtwo{\sqrt b-\sqrt c}{\tau}^2
 \leq\normone{b-c}{\tau}.
\]
\end{lemma}

\begin{proof}
Work in the tracial von Neumann algebra and put \(d=\sqrt b-\sqrt c\).
Traciality and \(b-c=\sqrt b\,d+d\sqrt c\) give
\[
 \tau(\sgn(d)(b-c))=\tau((\sqrt b+\sqrt c)|d|).
\]
Since \(\sqrt b+\sqrt c\geq d\) and
\(\sqrt b+\sqrt c\geq-d\), tracing against the positive and negative
parts \(d_+\) and \(d_-\), respectively, bounds the right-hand side
below by \(\tau(d_+^2)+\tau(d_-^2)=\tau(d^2)\).
Trace-norm duality bounds the left-hand side by
\(\normone{b-c}{\tau}\).
\end{proof}

\begin{proposition}[Weighted matrix blocks]\label{r:weighted:family}
Let \(K\subseteq\U(A)\) be finite and \(0<\eta<1\).
There are contractions \(v_1,\ldots,v_m\in A\) and nonnegative numbers
\(\lambda_1,\ldots,\lambda_m\) summing to one, chosen using only
\(K,\eta\), such that the map
\begin{equation}\label{r:eq:averaging-map}
 \Phi:D\longrightarrow D,\qquad
 \Phi(b)=\sum_{l=1}^m\lambda_l v_lbv_l^*,
\end{equation}
satisfies \eqref{r:eq:diagonal-mass} with \(\mu=\eta^2\).
For every projection \(r\in D\), there is a finite family consisting of
weights \(0<w_j\leq1\), nonzero projections \(q_j\leq r\), integers
\(k_j\geq1\), and rows \(U_j\in M_{1,k_j}(D)\), such that, with
\(p_j=U_jU_j^*\),
\begin{align}
 U_j^*U_j&=1_{k_j}\otimes q_j,\label{r:eq:block-isometry}\\
 \dist\big(U_j^*xU_j,M_{k_j}(\C)\otimes q_j\big)&<\eta
       &&(x\in K),\label{r:eq:block-compression}\\
 \normone{\sum_jw_jp_j-\Phi(r)}{\tau}&\leq\eta^2,
       \label{r:eq:block-mass}\\
 \sum_jw_j\normtwo{[x,p_j]}{\tau}^2&\leq6\eta
       &&(x\in K).\label{r:eq:block-energy}
\end{align}
The last two estimates hold for every \(\tau\in\T(D)\). The family may
be empty; the base projections \(q_j\) may repeat, and the projections
\(p_j\) need not be mutually orthogonal.
\end{proposition}

\begin{proof}
Put \(\mu=\eta^2\) and choose \(\Phi\) by
Lemma~\ref{r:weighted:diagonal}. In particular, the coefficient list is
fixed before \(r\). Choose \(0<\sigma<1\), then \(\beta>0\), so that
\begin{equation}\label{r:eq:block-parameters}
 \beta<\sigma/2,\qquad m(\sigma+2\beta)<\mu,\qquad
 \frac{2\beta}{\sigma}
       +\frac{(2+\sqrt2)\beta}{\sigma^2}<\eta.
\end{equation}
This is possible by choosing \(\sigma\) with \(m\sigma<\mu\) first.

\paragraph{Scalar matrix approximation.}
Partition \(r\) by Lemma~\ref{r:proj:compress}, using the finite list
\[
 v_l^*v_n,\quad v_l^*xv_n
 \qquad(1\leq l,n\leq m,\ x\in K).
\]
Take the compression tolerance smaller than \(\beta/(2m)\).
For a nonzero partition piece \(q\), set
\[
 V=[\sqrt{\lambda_l}\,v_lq]_{l=1}^m,\quad
 h_q=VV^*=\Phi(q),\quad M=V^*V,\quad N_x=V^*xV.
\]
These are contractions in their respective matrix spaces.
Apply any state of \(qDq\) entrywise to \(M,N_x\), obtaining scalar
matrices \(G,H_x\). Positivity and contractivity of a state at matrix
levels give \(0\leq G\leq1\) and \(\norm{H_x}\leq1\).
The compression tolerance implies
\begin{equation}\label{r:eq:scalar-gram}
 \norm{M-G\otimes q}<\beta,\qquad
 \norm{N_x-H_x\otimes q}<\beta.
\end{equation}
Indeed, each unweighted entry differs from its scalar state value by
less than twice the compression tolerance, and an \(m\)-by-\(m\) matrix
of entries of norm at most \(c\) has norm at most \(mc\).

\paragraph{Construction of the blocks.}
Rotate the columns of \(V\) by a scalar unitary diagonalizing \(G\).
Retain the columns corresponding to eigenvalues at least \(\sigma\),
forming a row \(W\) with \(k\) columns and scalar diagonal matrix \(G_+\).
The cutoff gives the uniform lower bound needed to normalize the retained
row by an inverse square root.
If \(k=0\), set \(b_q=0\). Otherwise put
\begin{equation}\label{r:eq:polar-block}
 X=W^*W\geq(\sigma/2)(1_k\otimes q),\qquad
 U=WX^{-1/2},\qquad b_q=U(G_+\otimes q)U^*.
\end{equation}
Then \(U^*U=1_k\otimes q\), so \(b_q\) is a finite-spectrum positive
contraction.
Writing \(W_-\) for the row of deleted columns, we have
\(\norm{W_-^*W_-}\leq\sigma+\beta\), and therefore
\[
 \tau(W_-W_-^*)=(\Tr\otimes\tau)(W_-^*W_-)
 \leq m(\sigma+\beta)\tau(q).
\]
The contribution is zero if no columns are deleted. Also
\[
 b_q-WW^*=U(G_+\otimes q-X)U^*
\]
has norm at most \(\beta\), supported on a projection of trace
\(k\tau(q)\). Consequently, including the case \(k=0\),
\begin{equation}\label{r:eq:finite-spectrum-close}
 \normone{b_q-h_q}{\tau}
 \leq m(\sigma+2\beta)\tau(q)\leq\mu\tau(q).
\end{equation}

We verify explicitly that the compression estimate is independent of
the corner and its matrix size. For positive invertible \(X,Y\geq a1\),
the resolvent formula for inverse square roots gives
\begin{equation}\label{r:eq:inverse-root}
 \norm{X^{-1/2}-Y^{-1/2}}
 \leq\frac{\norm{X-Y}}{\pi}
       \int_0^\infty\frac{t^{-1/2}}{(a+t)^2}\,dt
 =\frac{\norm{X-Y}}{2a^{3/2}}.
\end{equation}
Use \(Y=G_+\otimes q\) and \(a=\sigma/2\).
If \(H\otimes q\) denotes the corresponding retained corner of
\(H_x\otimes q\), then \(\norm H\leq1\), and
\eqref{r:eq:scalar-gram} gives
\(\norm{W^*xW-H\otimes q}<\beta\).
Expanding the difference between
\[
 U^*xU=X^{-1/2}(W^*xW)X^{-1/2}
 \quad\hbox{and}\quad Y^{-1/2}(H\otimes q)Y^{-1/2}
\]
and using \eqref{r:eq:inverse-root} bounds its norm by
\[
 \frac{2\beta}{\sigma}
 +\frac{(2+\sqrt2)\beta}{\sigma^2}<\eta.
\]
The second expression is a scalar matrix over \(q\). The tolerances
\(\sigma,\beta\) were chosen before the partition of \(r\), so the
compression bound is uniform over its pieces.

Write \(b_q\) as the finite sum of its nonzero upper spectral
projections, weighted by the lengths of their level intervals:
\begin{equation}\label{r:eq:layer-cake}
 b_q=\int_0^1 1_{(a,\infty)}(b_q)\,da
     =\sum_{j\in J_q}w_jp_j.
\end{equation}
Each \(p_j\) is the image of a scalar coordinate projection under \(U\).
It therefore has a subrow \(U_j\) satisfying
\eqref{r:eq:block-isometry}, with \(q_j=q\).
Taking matrix corners gives \eqref{r:eq:block-compression}.
Summing \eqref{r:eq:finite-spectrum-close} over the partition gives
\eqref{r:eq:block-mass}.

\paragraph{Averaged commutators.}
For a finite-spectrum
positive contraction \(b=\sum_s s e_s\), include its kernel projection
in this spectral resolution. Tracial orthogonality gives, for unitary
\(x\),
\begin{align}
 \int_0^1\normtwo{[x,1_{(a,\infty)}(b)]}{\tau}^2\,da
 &=\sum_{s,t}|s-t|\normtwo{e_sxe_t}{\tau}^2\notag\\
 &\leq2\sqrt{\tau(b)}\,\normtwo{[x,\sqrt b]}{\tau}\notag\\
 &\leq2\sqrt{\tau(b)\normone{b-xbx^*}{\tau}}.
 \label{r:eq:spectral-energy}
\end{align}
For the middle inequality, factor
\(|s-t|=|\sqrt s-\sqrt t|(\sqrt s+\sqrt t)\) and apply
Cauchy--Schwarz to the sum. The second squared sum is at most
\(2\sum_{s,t}(s+t)\normtwo{e_sxe_t}{\tau}^2=4\tau(b)\).
The last inequality is Lemma~\ref{r:weighted:square-root}.

For \(b=b_q\), equations \eqref{r:eq:diagonal-mass},
\eqref{r:eq:diagonal-central}, and \eqref{r:eq:finite-spectrum-close} give
\[
 \tau(b_q)\leq(1+\mu)\tau(q),\qquad
 \normone{b_q-xb_qx^*}{\tau}\leq3\mu\tau(q).
\]
Thus the contribution of \(J_q\) in \eqref{r:eq:block-energy} is at most
\(2\sqrt{3(1+\mu)}\,\eta\tau(q)\leq6\eta\tau(q)\).
Summing over \(q\) proves the assertion.
\end{proof}
